% !TEX root = ../main.tex
\chapter{Introduction to Spectroscopy}
\label{copy:chapter_spectroscopy}

\section{Spectroscopy and Historical Motivation}
\label{copy:sec:history}

\noindent
The origins of spectroscopy trace back to the observation of discrete spectral lines in atomic emission and absorption.
Bohr’s model (1913) provided the first quantitative interpretation by linking these lines to transitions between quantized energy levels, $\Delta E = h\nu$, thereby explaining the regularities of the Balmer and Rydberg series~\cite{bohr1913constitutionatomsmolecules, rydberg1890rybergsseriesrydberg1890xxxivstructurelinespectra, balmer1885notizuberspectrallinien}.
Although limited to hydrogen-like systems, this model established the idea that atomic structure could be inferred directly from the spectrum.

The emergence of wave mechanics soon gave spectroscopy a rigorous theoretical foundation.
In this framework, transition frequencies correspond to energy differences between eigenstates of the Hamiltonian, and selection rules follow from the symmetry properties of the transition operator~\cite{dirac1927quantumtheoryemission}. % TODO ?, wigner1927selectionrules, townes1955microwave}.
With advancing experimental resolution, subtle effects such as fine and hyperfine structure and, later, the Lamb shift became accessible, providing precise benchmarks that confirmed the predictions of quantum electrodynamics~\cite{lambretherford1947finestructurehydrogen, bethe1947electromagneticshiftenergy, hansch2006nobellecturepassion} and solidified the link between atomic spectra and fundamental theory.

Building on this foundation, the invention of the laser in 1960 introduced light sources of unprecedented coherence and spectral purity \todoref{...}%~\cite{maiman1960laser, townes1964laserprinciples}.
Lasers enabled Doppler-free and nonlinear techniques that overcame inhomogeneous broadening and opened the path to precision frequency metrology\todoref{...}%~\cite{hansch1975dopplerfree, hall2006nobel}.
When pulsed lasers became available, spectroscopy expanded from the frequency domain into the time domain, where ultrashort pulses allowed the direct observation of quantum coherence and energy-transfer dynamics\todoref{...}%~\cite{zewail1988femtochemistry, jonas2003twodimensionalfemtosecondspectroscopy}.
The development of photon-echo and pump–probe methods demonstrated that dephasing and coupling processes in condensed-phase systems could be resolved with femtosecond precision~\cite{mukamel1995principlesnonlinearoptical, cho2009twodimensionalopticalspectroscopy}.

These developments laid the foundation for coherent multidimensional spectroscopy, in which several time-ordered laser interactions generate nonlinear polarization signals that are Fourier-transformed along multiple time axes.
Two-dimensional electronic spectroscopy (2DES) is a central example: it correlates excitation and emission frequencies to disentangle overlapping transitions, reveal electronic couplings, and track coherence and spectral diffusion in complex molecular systems~\cite{jonas2003twodimensionalfemtosecondspectroscopy, hamm2005principlesnonlinearoptical, cho2009twodimensionalopticalspectroscopy, mukamel1995principlesnonlinearoptical, boyd2008contents}.
Such experiments probe the same dynamical information that can be accessed theoretically through open-quantum-system models, where system–bath interactions govern the loss of coherence and energy relaxation. \todoidea{shorten cause its not true}

In this thesis, these spectroscopic phenomena are analyzed theoretically using master-equation approaches that extend beyond simple perturbative treatments.
Particular focus is placed on Redfield-type theories, which describe relaxation and decoherence processes in systems coupled to thermal environments~\cite{redfield1957theory, breuer2002theoryopenquantum, may2011quantum}.
The derivation and physical interpretation of these equations are presented in the following chapter~\ref{chapter_open_quantum_systems}, which establishes the theoretical framework used for the simulations of two-dimensional electronic spectroscopy in hot and strongly coupled environments.


\section{Concept of Spectroscopy and Light--Matter Interaction}
\label{copy:sec:concept_light_matter}
\noindent
Throughout this thesis, spectroscopy means inferring structure and dynamics from the system’s optical polarization induced by external fields. We use a semiclassical model: a quantum system driven by a classical electric field. In Coulomb gauge and dipole approximation the Hamiltonian is
\begin{equation}
    H(t)=H_0 - \hat{\boldsymbol{\mu}}\cdot \boldsymbol{\varepsilon}(t),
    \label{copy:eq:H_dipole}
\end{equation}
where $H_0$ is the field-free Hamiltonian, $\hat{\boldsymbol{\mu}}$ the dipole operator, and $\boldsymbol{\varepsilon}(t)$ the electric field. The macroscopic polarization is the expectation value of the dipole,
\begin{equation}
    \boldsymbol{P}(t)=\mathrm{Tr}\!\left[\hat{\boldsymbol{\mu}}\rho(t)\right],
    \label{copy:eq:P_def}
\end{equation}
and, under the slowly varying envelope approximation, the detected signal field in a given phase-matched direction is proportional to the corresponding component of $\boldsymbol{P}(t)$ up to known prefactors \cite{mukamel1995principlesnonlinearoptical, boyd2008chapter6nonlinear}. We write the complex analytic field as $\boldsymbol{\varepsilon}(t)=\boldsymbol{\varepsilon}^{(+)}(t)+\boldsymbol{\varepsilon}^{(-)}(t)$ with $\boldsymbol{\varepsilon}^{(-)}=[\boldsymbol{\varepsilon}^{(+)}]^*$.

\section{Linear Spectroscopy: Response Functions and Lineshapes}
\label{copy:sec:linear}
\noindent
Expanding $\boldsymbol{P}$ in powers of the field gives
\begin{equation}
    \boldsymbol{P}(t)=\sum_{n\ge 1}\boldsymbol{P}^{(n)}(t),\qquad
    \boldsymbol{P}^{(n)}(t)=\varepsilon_0\int_0^\infty\!\!\cdots\!\int_0^\infty\!\!\mathrm{d}t_1\cdots\mathrm{d}t_n\;
    \boldsymbol{\chi}^{(n)}(t_1,\ldots,t_n)\,\prod_{k=1}^n \boldsymbol{\varepsilon}(t-\tau_k),
\end{equation}
with $\tau_k=\sum_{j=k}^n t_j$ and causal response tensors $\boldsymbol{\chi}^{(n)}$. For linear spectroscopy ($n=1$),
\begin{equation}
    P_i^{(1)}(t)=\varepsilon_0\int_0^\infty \mathrm{d}t_1\, \chi^{(1)}_{ij}(t_1)\,\varepsilon_j(t-t_1),
    \qquad
    \chi^{(1)}_{ij}(t)=\frac{i}{\hbar}\,\Theta(t)\,\langle[\hat{\mu}_i(t),\hat{\mu}_j(0)]\rangle,
    \label{copy:eq:chi1_comm}
\end{equation}
where $\langle\cdot\rangle$ denotes an equilibrium average with respect to $H_0$ and operators are in the interaction picture of $H_0$ \cite{mukamel1995principlesnonlinearoptical}. Fourier transforming gives
\begin{equation}
    \boldsymbol{P}^{(1)}(\omega)=\varepsilon_0\,\boldsymbol{\chi}^{(1)}(\omega)\,\boldsymbol{\varepsilon}(\omega),\quad
    \alpha(\omega)\propto \omega\,\Im\!\left[\mathrm{Tr}\,\boldsymbol{\chi}^{(1)}(\omega)\right],
\end{equation}
so the absorption spectrum is proportional to the imaginary part of the linear susceptibility.

\paragraph{Lineshape theory.}
For systems with diagonal system--bath coupling that is Gaussian and stationary, the linear absorption of a single transition with transition dipole $\mu$ can be written as
\begin{equation}
    I(\omega)\propto \Re\int_0^\infty \mathrm{d}t\; e^{i(\omega-\omega_0)t-g(t)-\Gamma t},
    \label{copy:eq:linear_lineshape}
\end{equation}
where $g(t)$ is the second-order cumulant (lineshape) function determined by the frequency--frequency correlation function, and $\Gamma$ accounts for lifetime broadening \cite{mukamel1995principlesnonlinearoptical, hamm2005principlesnonlinearoptical}. Equation~\eqref{eq:linear_lineshape} is exact for Gaussian fluctuations with linear coupling and is widely accurate when pure dephasing dominates. Inhomogeneous broadening enters by averaging over static detunings. The decomposition into diagonal (inhomogeneous) and antidiagonal (homogeneous) widths in 2D spectra follows from these same ingredients and is used below for interpretation.

\section{Extension to Nonlinear Spectroscopy: Third-Order Response and 2D Signals}
\label{copy:sec:third_order}
\noindent
Third-order spectroscopy probes $\boldsymbol{P}^{(3)}$, which is linear in the third-order response $\boldsymbol{R}^{(3)}$ and cubic in the fields. Using the standard response-function formalism,
\begin{align}
    P_i^{(3)}(t)                      & = \int_0^\infty\!\!\mathrm{d}t_3 \int_0^\infty\!\!\mathrm{d}t_2 \int_0^\infty\!\!\mathrm{d}t_1\;
    R^{(3)}_{i j k \ell}(t_3,t_2,t_1)\;
    \varepsilon_j(t-t_3)\,\varepsilon_k(t-t_3-t_2)\,\varepsilon_\ell(t-t_3-t_2-t_1), \label{copy:eq:P3_conv}                             \\
    R^{(3)}_{i j k \ell}(t_3,t_2,t_1) & =
    \left(-\frac{i}{\hbar}\right)^3
    \big\langle
    \big[\!\big[\!\big[\hat{\mu}_i(t_1{+}t_2{+}t_3),\hat{\mu}_j(t_2{+}t_3)\big],\hat{\mu}_k(t_3)\big],\hat{\mu}_\ell(0)\big]
    \big\rangle. \label{copy:eq:R3_def}
\end{align}
For impulsive, well-separated pulses one identifies four Liouville-space pathways for two-level optical transitions, commonly labeled $R_{1}\ldots R_{4}$ \cite{mukamel1995principlesnonlinearoptical, cho2009twodimensionalopticalspectroscopy}. In rotating-wave approximation the rephasing and nonrephasing combinations are selected by distinct phase-matching conditions:
\begin{align}
    \text{Rephasing:}\quad \boldsymbol{k}_S    & = -\boldsymbol{k}_1+\boldsymbol{k}_2+\boldsymbol{k}_3 \quad (R_2,R_3), \\
    \text{Nonrephasing:}\quad \boldsymbol{k}_S & = +\boldsymbol{k}_1-\boldsymbol{k}_2+\boldsymbol{k}_3 \quad (R_1,R_4).
\end{align}
Writing the three inter-pulse delays as $(t_1,t_2,t_3)\equiv(t_{\mathrm{coh}},T,t_{\mathrm{det}})$, the detected complex signal fields in each direction are proportional to the corresponding $\boldsymbol{P}^{(3)}$ component. Fourier transforms over $t_{\mathrm{coh}}$ and $t_{\mathrm{det}}$ yield the 2D spectra,
\begin{align}
    S_{R}(\omega_{\mathrm{coh}},T,\omega_{\mathrm{det}})
     & =\int \mathrm{d}t_{\mathrm{coh}}\,\mathrm{d}t_{\mathrm{det}}\;
    e^{+i\omega_{\mathrm{coh}} t_{\mathrm{coh}}}\,e^{-i\omega_{\mathrm{det}} t_{\mathrm{det}}}\;\varepsilon_{R}(t_{\mathrm{coh}},T,t_{\mathrm{det}}), \label{copy:eq:SR} \\
    S_{NR}(\omega_{\mathrm{coh}},T,\omega_{\mathrm{det}})
     & =\int \mathrm{d}t_{\mathrm{coh}}\,\mathrm{d}t_{\mathrm{det}}\;
    e^{-i\omega_{\mathrm{coh}} t_{\mathrm{coh}}}\,e^{-i\omega_{\mathrm{det}} t_{\mathrm{det}}}\;\varepsilon_{NR}(t_{\mathrm{coh}},T,t_{\mathrm{det}}), \label{copy:eq:SNR}
\end{align}
with opposite sign convention in the $t_{\mathrm{coh}}$ transform between $R$ and $NR$ branches \cite{cho2009twodimensionalopticalspectroscopy}. The absorptive spectrum combines both contributions,
\begin{equation}
    S_{\mathrm{abs}}(\omega_{\mathrm{coh}},T,\omega_{\mathrm{det}})=
    \Re\!\left\{S_R(\omega_{\mathrm{coh}},T,\omega_{\mathrm{det}})+S_{NR}(\omega_{\mathrm{coh}},T,\omega_{\mathrm{det}})\right\},
    \label{copy:eq:S_abs}
\end{equation}
which removes dispersive distortions and yields the quantity most directly compared to experiment \cite{jonas2003twodimensionalelectronicspectroscopy}.

\paragraph{Nonperturbative field simulation.}
In this thesis the optical \emph{field} is treated nonperturbatively by numerically propagating the master equation under the full time-dependent drive and then extracting $n$th-order signals by field-phase selection (phase cycling) and subtraction of lower orders. This gives access to the exact third-order signal in the chosen geometry without assuming impulsive $\delta$-pulses, while the \emph{system--bath} treatment is specified in the open-system chapter (Redfield and variants).

\section{Experimental Pulse Sequences and Phase Matching}
\label{copy:sec:phase_matching}
\noindent
Three pulses with carrier wavevectors $(\boldsymbol{k}_1,\boldsymbol{k}_2,\boldsymbol{k}_3)$ and controllable phases $(\phi_1,\phi_2,\phi_3)$ impinge on the sample. In a noncollinear ``boxcar’’ geometry, spatial selection enforces momentum conservation,
\begin{equation}
    \boldsymbol{k}_S=\pm \boldsymbol{k}_1 \pm \boldsymbol{k}_2 \pm \boldsymbol{k}_3,
    \label{copy:eq:PM_general}
\end{equation}
with the signs determined by whether a field acts on the ket or bra branch (RWA). In a collinear geometry, phase cycling replaces spatial selection: by repeating measurements with programmed phase sets and applying a discrete Fourier filter one isolates $(l,m,n)=(-1,+1,+1)$ for rephasing and $(+1,-1,+1)$ for nonrephasing pathways, respectively \cite{hamm2005principlesnonlinearoptical, cho2009twodimensionalopticalspectroscopy}. Heterodyne detection with a local oscillator of known phase yields complex-valued $S_R$ and $S_{NR}$, enabling the absorptive combination in Eq.~\eqref{eq:S_abs}.

\section{Interpretation of 2D Features: Coherences, Populations, Cross-Peaks}
\label{copy:sec:interpretation}
\noindent
2DES correlates the excitation frequency $\omega_{\mathrm{coh}}$ with the detection frequency $\omega_{\mathrm{det}}$ at fixed waiting time $T$.

\paragraph{Diagonal peaks.}
Features near $\omega_{\mathrm{det}}\!\approx\!\omega_{\mathrm{coh}}$ report the underlying linear absorption transitions. Their \emph{diagonal} width reflects static disorder (inhomogeneous broadening); their \emph{antidiagonal} width reflects homogeneous dephasing and lifetime broadening.

\paragraph{Cross-peaks and couplings.}
Off-diagonal peaks indicate coupling or population transfer between transitions. Static electronic coupling yields stationary cross-peaks whose phases and amplitudes depend on the relative transition dipoles and level structure. Time-dependent growth of cross-peaks with $T$ reveals relaxation or energy transfer pathways.

\paragraph{Coherences during $T$.}
Oscillatory modulation of peak amplitudes vs.\ $T$ signals coherent superpositions that survive into the waiting period (electronic or vibronic). The oscillation frequency matches level splittings in the active manifold. Damping of these oscillations quantifies decoherence rates.

\paragraph{Rephasing vs.\ nonrephasing contrast.}
Rephasing spectra (photon echo) are sensitive to inhomogeneous disorder and display echo narrowing along the antidiagonal; nonrephasing spectra emphasize homogeneous dynamics. The absorptive sum combines both for improved resolution and quantitative comparison.

\paragraph{Homogeneous vs.\ inhomogeneous decomposition.}
Under Gaussian fluctuations with linear coupling, the 2D lineshape can be expressed using the same lineshape function $g(t)$ entering Eq.~\eqref{eq:linear_lineshape}, explaining why antidiagonal widths report homogeneous dephasing while diagonal elongation reports static disorder \cite{mukamel1995principlesnonlinearoptical, hamm2005principlesnonlinearoptical}. This identification is exact for pure-dephasing Gaussian baths and remains a reliable guide more generally; deviations arise when spectral diffusion is non-Gaussian, when system--bath coupling is strong and nonperturbative, or when population lifetimes dominate.

\paragraph{Connection to open-system modeling.}
Decoherence and relaxation during $(t_1,T,t_3)$ depend on bath spectral densities and temperature. In Redfield-type theories the decay rates and frequency shifts are determined by system--bath couplings and bath correlation functions. At high temperature or for structured baths (e.g.\ Drude--Lorentz vs.\ Ohmic) the rates and spectral diffusion differ, altering antidiagonal widths, echo strength, and cross-peak kinetics. The open-system chapter specifies the generators and assumptions (Born--Markov, secular or not) used in the simulations; the present chapter provides the mapping from dynamics to observables.

\bigskip
\noindent\textbf{Scope and forward pointer.}
This chapter defines the optical observables and the pathway structure underlying 2DES. The following chapters specify (i) the electronic models used (two-level, four-level, and $N$-level one- and two-exciton manifolds), (ii) the open-system treatments (Redfield variants) with Drude--Lorentz and Ohmic baths, and (iii) the nonperturbative-in-field simulation protocol and numerical phase selection used to construct $S_R$, $S_{NR}$, and $S_{\mathrm{abs}}$ consistent with Eqs.~\eqref{eq:SR}--\eqref{eq:S_abs}.
